\documentclass{article}
\usepackage[utf8]{inputenc}
\usepackage[T1]{fontenc}
\usepackage{geometry}
\geometry{
    a4paper,
    left=2.5cm,
    right=2.5cm,
    top=2.5cm,
    bottom=2.5cm
}
\usepackage{graphicx}
\usepackage{float}
\usepackage{hyperref}
\usepackage[english]{babel}
\usepackage{amsmath}
\usepackage{amssymb}
\usepackage{array}
\usepackage{booktabs}
\usepackage{caption}
\usepackage{subcaption}
\usepackage{enumitem}
\usepackage{fancyhdr}
\usepackage{titlesec}
\usepackage{url}
\usepackage{textcomp}
\usepackage{xcolor}
\usepackage{listings}
\usepackage{multirow}
\usepackage{longtable}
\usepackage{tabularx}
\usepackage{times}

% Set up page style
\pagestyle{fancy}
\fancyhf{}
\rfoot{\thepage}
\renewcommand{\headrulewidth}{0pt}

% Title formatting
\titleformat{\chapter}[display]{\normalfont\bfseries}{\chaptertitle\thechapter}{0pt}{\Huge}
\titlespacing*{\chapter}{0pt}{50pt}{40pt}

% Document info
\title{Template Document}
\author{Academic Support Agent}
\date{\today}

\begin{document}
\maketitle

% Starten mit r mischer Nummerierung f r Formkapitel
\pagenumbering{Roman}

% =======================================================
% ===================== TITELSEITE ======================
% =======================================================
\begin{titlepage}
	\centering
	\vspace*{\fill}

	\includegraphics[width=0.7\textwidth]{logo_hsos.pdf}

	\vspace{2em}
	{\LARGE\textsc{\MYTITLE}\par}

	\vspace{2em}
	{
		\large
		Fakult t f r Ingenieurwissenschaften und Informatik \\
		der Hochschule Osnabr ck \\
		Studiengang \enquote{Verteilte und Mobile Anwendungen} \\
	}

	\vspace{2.5em}
	{\Large\textbf{(Master-/Bachelor-/Projekt-)arbeit}\par}

	\vspace{1em}
	% If Master/Bachelorarbeit
	zur Erlangung des akademischen Grades \\
	Master/Bachelor of Science/Engineering\par
	% else if Projektarbeit
	%	Im Modul \\
	%	--- hier Modulname ---\par

	\vspace{1.5em}
	vorgelegt von\par
	\vspace{1em}
	{\large\textbf{\MYAUTHOR}\par}

	\vspace*{\fill} % nach unten flie en
	\begin{center}
		\begin{tabular}{lll}
			\textbf{Matrikelnummer:} &  & XXXXXX                                \\
			\textbf{Ausgabedatum:}   &  & TT.MM.YYYY                            \\
			\textbf{Abgabedatum:}    &  & TT.MM.YYYY                            \\
			\textbf{Erstpr fer:}     &  & Name des 1. Pr fers inkl. akad. Titel \\
			\textbf{Zweitpr fer:}    &  & Name des 2. Pr fers inkl. akad. Titel \\
		\end{tabular}
	\end{center}
\end{titlepage}

% =======================================================
% ====================== DANKSAGUNG =====================
% =======================================================
\onehalfspacing						% Ab hier 1,5mm Abstand
\pagestyle{plain}					% Formatierung der Kopf und Fu zeilen
\setcounter{page}{1}				% Counter neu setzen, das ist die erste Seite

\addsec{Danksagungen}
\label{danksagungen}
% Generell optional, kann aber gerne eingef gt werden bei gr eren Arbeiten
% Bei Projektarbeiten sehr un blich
Ich danke allen unfassbar tollen Mitarbeitern im Labor f r Hochfrequenztechnik und Mobilkommunikation. Ihr seid die besten <3. Besonders m chte ich Daniel Brettschneider danken, der diese LaTeX Vorlage erstellt hat und Julian Dreyer, der diese Vorlage erweitert und an das modernere Design angepasst hat.
\pagebreak

% =======================================================
% ==================== ZUSAMMENFASSUNG ==================
% =======================================================

% Hier nicht mit sections arbeiten. Das macht unn tig Aufwand
\addsec{Zusammenfassung/Abstract}
Lorem ipsum dolor sit amet, consetetur sadipscing elitr, sed diam nonumy eirmod tempor invidunt ut labore et dolore magna aliquyam erat, sed diam voluptua. At vero eos et accusam et justo duo dolores et ea rebum. Stet clita kasd gubergren, no sea takimata sanctus est Lorem ipsum dolor sit amet. Lorem ipsum dolor sit amet, consetetur sadipscing elitr, sed diam nonumy eirmod tempor invidunt ut labore et dolore magna aliquyam erat, sed diam voluptua. At vero eos et accusam et justo duo dolores et ea rebum. Stet clita kasd gubergren, no sea takimata sanctus est Lorem ipsum dolor sit amet.

\minisec{Abstract}
Das Ganze auf Englisch.
\pagebreak

% Jetzt kann gerne eine Header Linie da sein
\renewcommand{\headrulewidth}{0.2pt}
% =======================================================
% ================== INHALTSVERZEICHNIS =================
% =======================================================
% Das Inhaltsverzeichnis auch im Inhaltsverzeichnis auftauchen lassen
\addcontentsline{toc}{section}{\texorpdfstring{Inhaltsverzeichnis}{Inhaltsverzeichnis}}
% Section Counter erh hen, wir haben ja eine neue section eingef gt
\addtocounter{section}{1}

% Maximal Tiefe 1 (max bis 1.1, nicht 1.1.1) im TOC
% ggf. auf 2  ndern, falls 1.1.1 gew nscht ist
\setcounter{tocdepth}{1}


% =======================================================
% ================== ABBILDUNGSVERZEICHNIS ==============
% =======================================================
% zu dem Inhaltsverzeichnis (TOC) hinzuf gen
\addcontentsline{toc}{chapter}{\listfigurename}

\pagestyle{fancy}

% =======================================================
% ================== TABELLENVERZEICHNIS ================
% =======================================================
\addcontentsline{toc}{chapter}{\listtablename}

\pagestyle{fancy}

% =======================================================
% ================== Codeverzeichnis ====================
% =======================================================
\renewcommand{\lstlistlistingname}{Quellcodeverzeichnis}
{\labelsep2cm\lstlistoflistings}
\addcontentsline{toc}{chapter}{\lstlistlistingname}


% =======================================================
% ================== Abk rzungsverzeichnis ==============
% =======================================================
\addchap{Abk rzungsverzeichnis}
\thispagestyle{plain}
\begin{acronym}[OSGi]						 % l ngste Abk rzung steht in eckigen Klammern
	\setlength{\itemsep}{-\parsep} 			 % geringerer Zeilenabstand
	\acro{IMU}{Inertial Measurment Unit}
 
	\acro{OSGi}{Open Service Gateway initiative}
	\acro{RSA}{Rivest Shamir Adleman}
\end{acronym}
% Nun kann mit den Befehlen \ac zB gearbeitet werden
\newpage

% =======================================================!
% =======================================================!
% =================== INHALTSBEGINN =====================!
% =======================================================!
% =======================================================!



% =======================================================
% ======================= EINLEITUNG ====================
% =======================================================
% Ab jetzt arabische Nummerierung. Zur cksetzen der Nummerierung
\setcounter{savepage}{\arabic{page}}
\pagenumbering{arabic}
\setcounter{page}{1}

% Ab hier nun stets den mystyle Pagestyle f r oben links
\pagestyle{mystyle}
\chapter{Einleitung}
\label{chap:einleitung}
Dieses Kapitel enth lt Beispiele zum Einf gen von Abbildungen, Tabellen, etc.

\section{Bilder}
\label{sec:bilder}
Zum Einf gen eines Bildes, siehe Abbildung \ref{fig:osgi}, wird die \textit{minipage}-Umgebung genutzt, da die Bilder so gut positioniert werden k nnen.

\vspace{1em}
\begin{minipage}{\linewidth}
	\centering
	\includegraphics[width=0.7\linewidth]{layering-osgi.png}
	\captionof{figure}[OSGi Architektur]{OSGi Architektur\footnotemark }
	\label{fig:osgi}
\end{minipage}
\footnotetext{Quelle: \url{http://www.osgi.org/Technology/WhatIsOSGi}}

\section{Tabellen}
\label{sec:tabellen}
In diesem Abschnitt wird eine Tabelle (siehe Tabelle \ref{tab:beispiel}) dargestellt.

\vspace{1em}
\begin{table}[!h]
	\centering
	\begin{tabular}{|l|l|l|}
		\hline
		\textbf{Name} & \textbf{Name} & \textbf{Name} \\
		\hline
		1             & 2             & 3             \\
		\hline
		4             & 5             & 6             \\
		\hline
		7             & 8             & 9             \\
		\hline
	\end{tabular}
	\caption{Beispieltabelle}
	\label{tab:beispiel}
\end{table}

\begin{table}[!h]
	\centering
	\begin{tabular}{|l|l|l|}
		\hline
		\textbf{Pin} & \textbf{Anschluss} & \textbf{Kabelfarbe} \\
		\hline
		Pin 1           & GND           & Wei           \\
		\hline
		Pin 2           & CAN L         & Braun         \\
		\hline
		Pin 3           & CAN H         & Blau          \\
        \hline
        Pin 4           & UBAT          & Schwarz       \\
		\hline
	\end{tabular}
	\caption{Pinout IMU}
	\label{tab:Pinout IMU}
\end{table}

\section{Auflistung}
\label{sec:auflistung}
F r Auflistungen wird die \textit{compactitem}-Umgebung genutzt, wodurch der Zeilenabstand zwischen den Punkten verringert wird.

\begin{compactitem}
	\item Nur
	\item ein
	\item Beispiel.
\end{compactitem}

\section{Listings}
\label{sec:listing}
Zuletzt ein Beispiel f r ein Listing, in dem Quellcode eingebunden werden kann, siehe Listing \ref{lst:arduino}.

\vspace{1em}
\begin{lstlisting}[caption=Arduino Beispielprogramm, label=lst:arduino]
int ledPin = 13;
void setup() {
    pinMode(ledPin, OUTPUT);
}
void loop() {
    digitalWrite(ledPin, HIGH);
    delay(500);
    digitalWrite(ledPin, LOW);
    delay(500);
}
\end{lstlisting}

\section{Tipps}
Die Quellen befinden sich in der Datei \textit{bibo.bib}. Eine Buch-, eine Paper- und eine Online-Quelle sind beispielhaft eingef gt. [Vgl. \cite{buch}, \cite{iotbds20}, \cite{online}]

Abk rzungen lassen sich nat rlich auch nutzen (\ac{OSGi}). Weiter oben im Latex-Code findet sich das Verzeichnis.

Mittels labels lassen sich auch Sektionen im Text referenzieren, z.B. \ref{sec:listing}, oder mit Autoref \autoref{sec:listing}.

% =======================================================
% ====================== Grundlagen =====================
% =======================================================
\chapter{Grundlagen}
\label{chap:grundlagen}
Lorem ipsum dolor sit amet.

\section{Unterkapitel}
\label{sec:u1}
Lorem ipsum dolor sit amet, consetetur sadipscing elitr, sed diam nonumy eirmod tempor invidunt ut labore et dolore magna aliquyam erat, sed diam voluptua. At vero eos et accusam et justo duo dolores et ea rebum. Stet clita kasd gubergren, no sea takimata sanctus est Lorem ipsum dolor sit amet. Lorem ipsum dolor sit amet, consetetur sadipscing elitr, sed diam nonumy eirmod tempor invidunt ut labore et dolore magna aliquyam erat, sed diam voluptua. At vero eos et accusam et justo duo dolores et ea rebum. Stet clita kasd gubergren, no sea takimata sanctus est Lorem ipsum dolor sit amet.

\subsection{Unterunterkapitel}
\label{sec:u2}
Lorem ipsum dolor sit amet, consetetur sadipscing elitr, sed diam nonumy eirmod tempor invidunt ut labore et dolore magna aliquyam erat, sed diam voluptua. At vero eos et accusam et justo duo dolores et ea rebum. Stet clita kasd gubergren, no sea takimata sanctus est Lorem ipsum dolor sit amet. Lorem ipsum dolor sit amet, consetetur sadipscing elitr, sed diam nonumy eirmod tempor invidunt ut labore et dolore magna aliquyam erat, sed diam voluptua. At vero eos et accusam et justo duo dolores et ea rebum. Stet clita kasd gubergren, no sea takimata sanctus est Lorem ipsum dolor sit amet.


% =======================================================
% ====================== Durchf hrung =====================
% =======================================================
\chapter{Durchf hrung}
\label{chap:durchf hrung}
Lorem ipsum dolor sit amet.

\section{Unterkapitel}
\label{sec:u3}
Lorem ipsum dolor sit amet, consetetur sadipscing elitr, sed diam nonumy eirmod tempor invidunt ut labore et dolore magna aliquyam erat, sed diam voluptua. At vero eos et accusam et justo duo dolores et ea rebum. Stet clita kasd gubergren, no sea takimata sanctus est Lorem ipsum dolor sit amet. Lorem ipsum dolor sit amet, consetetur sadipscing elitr, sed diam nonumy eirmod tempor invidunt ut labore et dolore magna aliquyam erat, sed diam voluptua. At vero eos et accusam et justo duo dolores et ea rebum. Stet clita kasd gubergren, no sea takimata sanctus est Lorem ipsum dolor sit amet.

\subsection{Unterkapitel}
\label{subsec:uu1}
Lorem ipsum dolor sit amet, consetetur sadipscing elitr, sed diam nonumy eirmod tempor invidunt ut labore et dolore magna aliquyam erat, sed diam voluptua. At vero eos et accusam et justo duo dolores et ea rebum. Stet clita kasd gubergren, no sea takimata sanctus est Lorem ipsum dolor sit amet. Lorem ipsum dolor sit amet, consetetur sadipscing elitr, sed diam nonumy eirmod tempor invidunt ut labore et dolore magna aliquyam erat, sed diam voluptua. At vero eos et accusam et justo duo dolores et ea rebum. Stet clita kasd gubergren, no sea takimata sanctus est Lorem ipsum dolor sit amet.

% Kommt nicht im TOC vor!
\subsubsection{Unterkapitel}
\label{subsubsec:uuu1}
Lorem ipsum dolor sit amet, consetetur sadipscing elitr, sed diam nonumy eirmod tempor invidunt ut labore et dolore magna aliquyam erat, sed diam voluptua. At vero eos et accusam et justo duo dolores et ea rebum. Stet clita kasd gubergren, no sea takimata sanctus est Lorem ipsum dolor sit amet. Lorem ipsum dolor sit amet, consetetur sadipscing elitr, sed diam nonumy eirmod tempor invidunt ut labore et dolore magna aliquyam erat, sed diam voluptua. At vero eos et accusam et justo duo dolores et ea rebum. Stet clita kasd gubergren, no sea takimata sanctus est Lorem ipsum dolor sit amet.


% =======================================================
% ====================== Analyse ========================
% =======================================================
\chapter{Analyse}
\label{chap:analyse}
Lorem ipsum dolor sit amet.

\section{Unterkapitel}
\label{sec:u4}
Lorem ipsum dolor sit amet, consetetur sadipscing elitr, sed diam nonumy eirmod tempor invidunt ut labore et dolore magna aliquyam erat, sed diam voluptua. At vero eos et accusam et justo duo dolores et ea rebum. Stet clita kasd gubergren, no sea takimata sanctus est Lorem ipsum dolor sit amet. Lorem ipsum dolor sit amet, consetetur sadipscing elitr, sed diam nonumy eirmod tempor invidunt ut labore et dolore magna aliquyam erat, sed diam voluptua. At vero eos et accusam et justo duo dolores et ea rebum. Stet clita kasd gubergren, no sea takimata sanctus est Lorem ipsum dolor sit amet.

\subsection{Unterunterkapitel}
\label{subsec:uu2}
Lorem ipsum dolor sit amet, consetetur sadipscing elitr, sed diam nonumy eirmod tempor invidunt ut labore et dolore magna aliquyam erat, sed diam voluptua. At vero eos et accusam et justo duo dolores et ea rebum. Stet clita kasd gubergren, no sea takimata sanctus est Lorem ipsum dolor sit amet. Lorem ipsum dolor sit amet, consetetur sadipscing elitr, sed diam nonumy eirmod tempor invidunt ut labore et dolore magna aliquyam erat, sed diam voluptua. At vero eos et accusam et justo duo dolores et ea rebum. Stet clita kasd gubergren, no sea takimata sanctus est Lorem ipsum dolor sit amet.


% =======================================================
% ======================== Fazit ========================
% =======================================================
\chapter{Fazit}
\label{chap:fazit}
Lorem ipsum dolor sit amet.

\section{Unterkapitel}
\label{sec:u5}
Lorem ipsum dolor sit amet, consetetur sadipscing elitr, sed diam nonumy eirmod tempor invidunt ut labore et dolore magna aliquyam erat, sed diam voluptua. At vero eos et accusam et justo duo dolores et ea rebum. Stet clita kasd gubergren, no sea takimata sanctus est Lorem ipsum dolor sit amet. Lorem ipsum dolor sit amet, consetetur sadipscing elitr, sed diam nonumy eirmod tempor invidunt ut labore et dolore magna aliquyam erat, sed diam voluptua. At vero eos et accusam et justo duo dolores et ea rebum. Stet clita kasd gubergren, no sea takimata sanctus est Lorem ipsum dolor sit amet.

\subsection{Unterunterkapitel}
\label{sec:uu3}
Lorem ipsum dolor sit amet, consetetur sadipscing elitr, sed diam nonumy eirmod tempor invidunt ut labore et dolore magna aliquyam erat, sed diam voluptua. At vero eos et accusam et justo duo dolores et ea rebum. Stet clita kasd gubergren, no sea takimata sanctus est Lorem ipsum dolor sit amet. Lorem ipsum dolor sit amet, consetetur sadipscing elitr, sed diam nonumy eirmod tempor invidunt ut labore et dolore magna aliquyam erat, sed diam voluptua. At vero eos et accusam et justo duo dolores et ea rebum. Stet clita kasd gubergren, no sea takimata sanctus est Lorem ipsum dolor sit amet.


% =======================================================
% ===================== Literatur =======================
% =======================================================
% Ab hier nun wieder r mische Seitenzahlen
\newpage
\pagenumbering{Roman}
\setcounter{page}{\thesavepage}


% Wenn nicht anders gew nscht, wird hier IEEE Citation verwendet




% =======================================================
% ======================= Anhang ========================
% =======================================================

\newpage
\addchap{Anhang}
\label{anhang}

\addsec{Inhalt der CD}
Hier steht die Inhaltsstruktur der beigelegten CD. Im Folgenden ist ein beispielhafter Aufbau dazu dargestellt:

\begin{table}[h]
	\begin{tabular}{|p{0.3\linewidth}|p{0.3\linewidth}|p{0.4\linewidth}|}
		%\begin{tabular}{|l|l|l|}
		\hline
		Ordnerverzeichnis                               & Dateien                                & Beschreibung                                         \\ \hline
		\multirow{2}{*}{\textbackslash{}Bachelorarbeit} & Bachelorarbeit\_Name.pdf               & Die Bachelorarbeit im Portable Document Format (PDF) \\
		                                                & \textbackslash{}Bachelorarbeit\_Source & Die LaTeX-Quelldateien der Bachelorarbeit            \\ \hline
		\textbackslash{}Referenzen                      & *.pdf                                  & Beschaffte Texte im Portable Document Format (PDF)   \\ \hline
		\multirow{2}{*}{\textbackslash{}Source\_Code}   & \textbackslash{}programm\_1            & Programm 1                                           \\
		                                                & \textbackslash{}generated\_data        & Repository f r generierte Daten                      \\\hline
	\end{tabular}
	\label{tab:anhang-inhalt-cd}
	\caption{Inhalt der CD}
\end{table}


\addsec{Weitere Anh nge}

% =======================================================
% ===================== Erkl rung =======================
% =======================================================
\newpage
\addchap{Eidesstattliche Erkl rung}
\label{erklaerung}

Hiermit erkl re ich an Eides statt, dass ich die vorliegende Arbeit selbst ndig und ohne fremde Hilfe angefertigt habe. Die aus fremden Quellen direkt oder indirekt  bernommenen Gedanken sind als solche einzeln kenntlich gemacht. Es wurden keine anderen als die angegebenen Quellen und Hilfsmittel benutzt. Die Arbeit wurde bisher keiner anderen Pr fungsbeh rde vorgelegt und auch nicht ver ffentlicht.
\vspace*{2em}

Osnabr ck, \hrulefill\qquad Unterschrift: \hrulefill
\vspace*{3em}


% Optional, muss nicht
\section*{Urheberrechtliche Einwilligungserkl rung}

Hiermit erkl re ich, dass ich damit einverstanden bin, dass meine Arbeit zum Zwecke des Plagiatsschutzes bei der Fa. Ephorus BV bis zu 5 Jahren in einer Datenbank f r die Hochschule Osnabr ck archiviert werden kann. Diese Einwilligung kann jederzeit widerrufen werden.
\vspace*{2em}

Osnabr ck, \hrulefill\qquad Unterschrift: \hrulefill
\end{document}